\section*{Capítulo \ref{ch:b3:lagrangian-mechanics} --- Mecânica lagrangiana}

\begin{solution}{exo:b3:lagrangian-mechanics:1}
(a) 2 (dois ângulos sobre a superfície da tigela). (b) 1: a abscissa
$x$ ao longo da encosta, com o ângulo de rotação preso a ela pelo
rolamento, $x =
R\phi$. (c) 3: $X$, $\theta_1$, $\theta_2$. (d) 5: três para o centro,
dois para a direção da haste. (e) 2: um ângulo cada. A condição de
rolamento é uma relação entre \emph{velocidades}, $\dot x =
R\dot\phi$; neste problema plano ela se integra de imediato em $x = R\phi +
\text{const}$, portanto é holônoma. (Para uma bola que rola sobre um
plano ela não se integra, e a mecânica \omterm{def:b3:lagrangian-mechanics:action}{lagrangiana} precisa de uma
extensão.)
\end{solution}

\begin{solution}{exo:b3:lagrangian-mechanics:2}
(a) $[L] = \unit{J}$; $[p_\theta] = \unit{kg.m^2/s}$ --- um momento
angular, pois $\theta$ é adimensional. (b) $p_\theta =
m\ell^2\dot\theta$ é o momento angular da massa em torno do pivô. (c)
$\ddot\theta = -(g/\ell)\theta$ dá $T = 2\pi\sqrt{\ell/g}$. (d) Com
$E_p$ medida a partir do equilíbrio, uma oscilação harmônica tem médias
temporais iguais de energia cinética e potencial, de modo que $\langle
L\rangle = 0$ e $S = 0$ ao longo de um período inteiro --- o cálculo
confirma: $S = \int_0^T\tfrac12 m\ell^2\theta_0^2[\omega^2\sin^2
\omega t - (g/\ell)\cos^2\omega t]\,\dd t = 0$, já que $\omega^2 =
g/\ell$.
\end{solution}

\begin{solution}{exo:b3:lagrangian-mechanics:3}
(a) Seja $x$ a descida de $m_1$ (de modo que $m_2$ sobe de $x$): $L =
\tfrac12(m_1 + m_2)\dot x^2 + (m_1 - m_2)gx$. (b) $(m_1 + m_2)\ddot x =
(m_1 - m_2)g$: $a = (m_1 - m_2)g/(m_1 + m_2)$. (c) A tração age nas
duas extremidades de um fio inextensível: em qualquer deslocamento
permitido os seus trabalhos se cancelam, e por isso ela nunca entra em
$L$. (d) Newton aplicado só a $m_1$: $T =
m_1(g - a) = 2m_1m_2g/(m_1 + m_2)$.
\end{solution}

\begin{solution}{exo:b3:lagrangian-mechanics:4}
(a) $\varphi$ e $z$: $p_\varphi = mr^2\dot\varphi$ (o momento angular
em torno do eixo) e $p_z = m\dot z$ se conservam. (b) $r$ aparece em
$L$ através de $r^2\dot\varphi^2$, logo não é cíclica: $\dot p_r =
mr\dot\varphi^2 \neq 0$. Nada está errado --- $p_r = m\dot r$ é a
componente radial de um vetor constante $\vect p$, e uma componente ao
longo de uma direção que \emph{gira} não precisa ser constante. (c) $m\ddot r =
mr\dot\varphi^2$: o termo centrífugo, o preço de usar coordenadas
presas a direções que giram. (d) Para uma reta à distância $b$
percorrida com velocidade $v$: $r = \sqrt{b^2 + v^2t^2}$, $r^2\dot\varphi = bv$;
então $\ddot r = b^2v^2/r^3 = r(bv/r^2)^2 = r\dot\varphi^2$.
\end{solution}

\begin{solution}{exo:b3:lagrangian-mechanics:5}
(a) Posição do bloco $(X + s\cos\alpha,\ -s\sin\alpha)$, logo
\[
L = \tfrac12 M\dot X^2 + \tfrac12 m\big(\dot X^2 +
2\dot X\dot s\cos\alpha + \dot s^2\big) + mgs\sin\alpha .
\]
(b) $X$ é cíclica: $P = (M + m)\dot X + m\dot s\cos\alpha$ se conserva
--- o momento horizontal total, pois nenhuma força horizontal externa
atua. (c) A equação de $s$ é $\ddot s + \ddot
X\cos\alpha = g\sin\alpha$; com $\ddot X = -m\ddot s\cos\alpha/(M +
m)$ vindo de (b),
\[
\ddot s = \frac{(M + m)g\sin\alpha}{M + m\sin^2\alpha} , \qquad
\ddot X = -\frac{mg\sin\alpha\cos\alpha}{M + m\sin^2\alpha} .
\]
(d) $M \to \infty$: $\ddot s \to g\sin\alpha$, o plano inclinado fixo;
$\alpha \to 90^\circ$: $\ddot s \to g$, $\ddot X \to 0$ --- queda livre
ao longo de uma face vertical, sem empurrar a cunha.
\end{solution}

\begin{solution}{exo:b3:lagrangian-mechanics:6}
(a) $L = \tfrac12 m\ell^2(\dot\theta^2 + \sin^2\theta\,\dot\varphi^2) +
mg\ell\cos\theta$. (b) $\varphi$ é cíclica: $p_\varphi =
m\ell^2\sin^2\theta\,\dot\varphi$, o momento angular vertical. (c)
Eliminando $\dot\varphi$, $\tfrac12 m\ell^2\dot\theta^2 +
U_{\text{eff}}(\theta)$ se conserva, com
\[
U_{\text{eff}}(\theta) = \frac{p_\varphi^2}{2m\ell^2\sin^2\theta}
 - mg\ell\cos\theta :
\]
uma barreira em $\theta = 0$ e $\theta = \pi$ (para $p_\varphi \neq 0$)
com um mínimo entre elas --- a massa nuta entre dois círculos. (d) No
mínimo, $\dot\theta = 0$ com $\theta = \theta_0$ constante:
$p_\varphi^2\cos\theta_0/m\ell^2\sin^3\theta_0 = mg\ell\sin\theta_0$;
inserindo $p_\varphi = m\ell^2\sin^2\theta_0\,\omega$ obtém-se
$\ell\omega^2\cos\theta_0 = g$.
\end{solution}

\begin{solution}{exo:b3:lagrangian-mechanics:7}
(a) Massa em $(X + \ell\sin\theta,\ -\ell\cos\theta)$:
\[
L = \tfrac12(M + m)\dot X^2 + m\ell\cos\theta\,\dot X\dot\theta
 + \tfrac12 m\ell^2\dot\theta^2 + mg\ell\cos\theta .
\]
(b) $X$ é cíclica: $(M + m)\dot X + m\ell\cos\theta\,\dot\theta$ se
conserva --- o momento horizontal do sistema inteiro (o trilho só
empurra verticalmente). (c) Pequenos ângulos: $(M + m)\ddot X +
m\ell\ddot\theta = 0$ e $\ell\ddot\theta + \ddot X + g\theta = 0$;
eliminando $\ddot X$, $\ell\ddot\theta\,[1 - m/(M + m)] = -g\theta$,
logo $\Omega^2 = (M + m)g/M\ell = (1 + m/M)\,g/\ell$. (d) $M \to \infty$:
o pivô fixo, $\Omega^2 = g/\ell$. Para $M$ pequeno, o carrinho recua em
sentido oposto ao da massa e a oscilação se dá em torno de um ponto
entre os dois (o centro de massa fixo), o que encurta o pêndulo efetivo
--- daí a frequência mais alta, que diverge quando $M \to 0$.
\end{solution}

\begin{solution}{exo:b3:lagrangian-mechanics:8}
(a) $\vect A = \tfrac12 B(-y, x, 0)$:
$L = \tfrac12 m(\dot x^2 + \dot y^2 + \dot z^2) + \tfrac12 qB(x\dot y -
y\dot x)$. (b) A equação de $x$: $\dd(m\dot x - \tfrac12 qBy)/\dd t =
\tfrac12 qB\dot y$, isto é, $m\ddot x = qB\dot y$; do mesmo modo $m\ddot y =
-qB\dot x$: movimento circular com $\omega_{\text{c}} = qB/m$ e $\dot z$
constante. (c) $p_x = m\dot x - \tfrac12 qBy \neq m\dot x$; nem $x$
nem $y$ é cíclica ($\partial L/\partial x = \tfrac12 qB\dot y$), de
modo que nenhum dos momentos se conserva --- só combinações como $m\dot x -
qBy$ (verifique: sua derivada se anula). (d) Em coordenadas
cilíndricas, $A_\varphi = \tfrac12 Br$:
$L = \tfrac12 m(\dot r^2 + r^2\dot\varphi^2 + \dot z^2) + \tfrac12
qBr^2\dot\varphi$; $\varphi$ é cíclica e $p_\varphi = mr^2\dot\varphi +
\tfrac12 qBr^2$. Num círculo de raio $R$ centrado no eixo,
$\dot\varphi = -\omega_{\text{c}}$, logo $p_\varphi = -qBR^2 + \tfrac12
qBR^2 = -\tfrac12 qBR^2$.
\end{solution}

\begin{solution}{exo:b3:lagrangian-mechanics:9}
(a) $h = \dot\theta\,\partial L/\partial\dot\theta - L = \tfrac12
mR^2\dot\theta^2 - \tfrac12 m\omega^2R^2\sin^2\theta - mgR\cos\theta =
\tfrac12 mR^2\dot\theta^2 + U_{\text{eff}}$; $\dd h/\dd t =
\dot\theta\,[mR^2\ddot\theta + U_{\text{eff}}'(\theta)] = 0$ pela
equação de movimento. (b) $E = \tfrac12 mR^2\dot\theta^2 + \tfrac12
m\omega^2R^2\sin^2\theta - mgR\cos\theta = h +
m\omega^2R^2\sin^2\theta$. (c) $P = \dd E/\dd t =
m\omega^2R^2\sin2\theta\,\dot\theta$: positiva enquanto a conta se
afasta do eixo (o motor trabalha contra a inércia da conta), negativa
na volta. (d) $U_{\text{eff}}'' = mgR\cos\theta -
m\omega^2R^2\cos2\theta$; em $\cos\theta_{\text{eq}} = g/\omega^2R$
isso vale $m\omega^2R^2\sin^2\theta_{\text{eq}}$, logo
$\omega_{\text{osc}} = \omega\sin\theta_{\text{eq}} = \omega\sqrt{1 -
(g/\omega^2R)^2} \to 0$ quando $\omega^2 \to g/R$: a força restauradora
se achata exatamente quando os poços se fundem.
\end{solution}

\begin{solution}{exo:b3:lagrangian-mechanics:10}
(a) Posições $x_2 = \ell(\sin\theta_1 + \sin\theta_2)$, etc.; guardando
os termos quadráticos, o termo cruzado de velocidade é
$m\ell^2\dot\theta_1\dot\theta_2$, o que dá a $L$ enunciada. (b)
$2\ddot\theta_1 + \ddot\theta_2 = -2\omega_0^2\theta_1$ e
$\ddot\theta_1 + \ddot\theta_2 = -\omega_0^2\theta_2$, com $\omega_0^2
= g/\ell$. (c) Inserindo $\theta_i = a_i\cos\Omega t$:
$\det\begin{pmatrix} 2\omega_0^2 - 2\Omega^2 & -\Omega^2\\ -\Omega^2 &
\omega_0^2 - \Omega^2\end{pmatrix} = \Omega^4 - 4\omega_0^2\Omega^2 +
2\omega_0^4 = 0$, logo $\Omega_\pm^2 = (2 \mp \sqrt2)\omega_0^2$; a
segunda linha dá $a_2/a_1 = \Omega^2/(\omega_0^2 - \Omega^2) =
\pm\sqrt2$: massas juntas (modo lento), massas opostas (modo rápido).
(d) Em grande amplitude, os acoplamentos em $\sin$ e $\cos$ tornam as
equações não lineares; as soluções deixam de se superpor e condições
iniciais vizinhas se afastam exponencialmente (caos). A energia
continua exatamente conservada --- $L$ não tem tempo explícito --- o
caos diz respeito à previsibilidade, não à conservação.
\end{solution}

\begin{solution}{exo:b3:lagrangian-mechanics:11}
(a) $v = \sqrt{2gy}$ e $\dd s = \sqrt{1 + y'^2}\,\dd x$ dão o
funcional. (b) O cálculo da
\cref{prop:b3:lagrangian-mechanics:energy} com $x$ no lugar do tempo: $\dd
h/\dd x = -\partial F/\partial x = 0$. (c) $h = -1/\sqrt{2gy(1 +
y'^2)}$, logo $y(1 + y'^2) = 2r$. Para a cicloide, $y' =
\sin\phi/(1 - \cos\phi)$ e $1 + y'^2 = 2/(1 - \cos\phi)$, de onde
$y(1 + y'^2) = 2r$. (d) $\dd t = \dd s/v = \sqrt{r/g}\,\dd\phi$ (toda a
dependência em $\phi$ se cancela), de modo que o fundo ($\phi = \pi$) é
alcançado em $\pi\sqrt{r/g}$ --- qualquer que seja o ponto de partida:
a cicloide é também a tautócrona. A rampa reta até $(\pi r, 2r)$ leva $\sqrt{\pi^2
+ 4}\,\sqrt{r/g} \approx 3.72\sqrt{r/g}$, cerca de $18\%$ a mais que
$\pi\sqrt{r/g}$.
\end{solution}

\begin{solution}{exo:b3:lagrangian-mechanics:12}
(a) $\dd t = \dd s/(c/n)$. (b) $h = -n(y)/\sqrt{1 + y'^2}$ se
conserva; $1/\sqrt{1 + y'^2} = \cos\theta$ ($\theta$ é o ângulo de
inclinação) $= \sin i$ com $i$ medido da vertical: $n\sin i = \text{const}$ ---
a lei de Snell, aplicada de modo contínuo. (c) Para um raio quase
horizontal, $n\cos\theta \approx \text{const}$ com $\theta$ pequeno dá
$\theta\,\dd\theta = \dd n/n \approx \beta\,\dd y$; como $\dd y =
\theta\,\dd x$, a curvatura é $\dd\theta/\dd x = \beta$, isto é, $R =
1/\beta \approx \qty{83}{km}$, encurvando-se para \emph{cima} (rumo aos
$n$ maiores). (d) Um raio que parte do olho e tangencia a estrada após
encurvar-se com raio $R$ toca-a em $d = \sqrt{2hR} = \sqrt{2 \times 1.2 \times
8.3e4} \approx \qty{450}{m}$: além dessa distância a própria estrada
não é vista --- vê-se o céu refratado para cima, a ``água'' cintilante.
\end{solution}

\begin{solution}{pb:b3:lagrangian-mechanics:1}
\textbf{1.} $L = \tfrac12 m\ell^2\dot\theta^2 + mg\ell\cos\theta$;
$\ddot\theta = -(g/\ell)\sin\theta$.
\textbf{2.} $f_0 = \sqrt{g/\ell}/2\pi = \qty{0.79}{Hz}$.
\textbf{3.} A suspensão é fixa, logo $h = E = \tfrac12
m\ell^2\dot\theta^2 - mg\ell\cos\theta$; do fundo ao topo, $\tfrac12
mv^2 = 2mg\ell$: $v = 2\sqrt{g\ell} = \qty{4.0}{m/s}$.
\textbf{4.} $\ddot\epsilon = +(g/\ell)\epsilon$: $\epsilon \propto
\eu^{t/\tau}$ com $\tau = \sqrt{\ell/g} = \qty{0.20}{s}$.
\textbf{5.} O equilíbrio invertido existe, mas é exponencialmente
instável: qualquer inclinação, por menor que seja, multiplica-se por
$e$ a cada \qty{0.2}{s}.
\textbf{6.} Massa em $(\ell\sin\theta,\ a\cos\Omega t - \ell\cos\theta)$;
derive e eleve ao quadrado: o $\vect v^{\,2}$ enunciado, com o termo
cruzado $-2a\Omega\ell\sin(\Omega t)\sin\theta\,\dot\theta$.
\textbf{7.} Se $L' = L + \dd F(q,t)/\dd t$, a \omterm{def:b3:lagrangian-mechanics:action}{ação} muda de
$F(q_2, t_2) - F(q_1, t_1)$, constante sob variações com extremidades
fixas: os mesmos caminhos estacionários, as mesmas equações.
\textbf{8.} O termo cruzado vale $-ma\Omega\ell\sin(\Omega t)\sin\theta\,
\dot\theta = \dd[ma\Omega\ell\sin(\Omega t)\cos\theta]/\dd t -
ma\Omega^2\ell\cos(\Omega t)\cos\theta$; descartando a derivada total
e os termos que só dependem de $t$ ($\tfrac12 ma^2\Omega^2\sin^2\Omega t$ e
$-mga\cos\Omega t$), resta a $L$ enunciada.
\textbf{9.} $m\ell^2\ddot\theta = -mg\ell\sin\theta +
ma\Omega^2\ell\cos(\Omega t)\sin\theta$. Como $\ddot y_{\text{s}} =
-a\Omega^2\cos\Omega t$, isso é $\ddot\theta =
-[(g + \ddot y_{\text{s}})/\ell]\sin\theta$: no referencial do pivô, a
gravidade aparente oscila.
\textbf{10.} $L$ passa a depender explicitamente de $t$: $h$ não se
conserva, nem $E$ --- o agitador injeta e retira energia pelo pivô.
\textbf{11.} $a/\ell = 0.05$; $\Omega = \qty{251}{rad/s}$ contra
$\omega_0 = \qty{4.9}{rad/s}$: razão $51$. Num período de excitação
(\qty{25}{ms}) a gravidade mal altera $\dot\theta$: a massa é lenta
demais para acompanhar e apenas estremece.
\textbf{12.} $\ddot\xi$ é a maior derivada ($\propto\Omega^2$) e a
excitação é o maior termo de força: $\ddot\xi =
(a\Omega^2/\ell)\cos(\Omega t)\sin\Theta$.
\textbf{13.} Integrando duas vezes com $\Theta$ fixo: $\xi =
-(a/\ell)\cos(\Omega t)\sin\Theta$, de amplitude no máximo $a/\ell =
0.05$: um tremor de dois graus.
\textbf{14.} $\sin\theta \approx \sin\Theta + \xi\cos\Theta$, logo
\[
\ddot\Theta + \ddot\xi = -\frac{g}{\ell}(\sin\Theta + \xi\cos\Theta)
 + \frac{a\Omega^2}{\ell}\cos\Omega t\,(\sin\Theta + \xi\cos\Theta) .
\]
\textbf{15.} A média elimina $\ddot\xi$, $\langle\cos\Omega t\rangle$
e $\langle\xi\rangle$; o termo cruzado que sobrevive é
$(a\Omega^2/\ell)\cos\Theta\,\langle\xi\cos\Omega t\rangle =
-(a^2\Omega^2/2\ell^2)\sin\Theta\cos\Theta$, o que dá a equação
enunciada para $\Theta$.
\textbf{16.} $\ddot\Theta = -(1/m\ell^2)\,U_{\text{eff}}'(\Theta)$ com
$U_{\text{eff}} = mg\ell[-\cos\Theta + (a^2\Omega^2/4g\ell)
\sin^2\Theta]$ --- derive para verificar.
\textbf{17.} O mesmo termo em $\sin^2$ do aro, mas com o sinal
\emph{oposto}: a rotação cavava poços nos flancos e só podia achatar o
fundo; a sacudida vertical enrijece o poço do fundo e cava um novo poço
no \emph{topo}.
\textbf{18.} Excitação lenta ($a^2\Omega^2 < 2g\ell$): mínimo em $\Theta =
0$, máximo em $\pi$ --- nada de novo. Excitação rápida: mínimos em $0$
\emph{e} em $\pi$, separados por máximos em $\cos\Theta^* =
-2g\ell/a^2\Omega^2$; tanto o pêndulo pendente quanto o pêndulo em pé
oscilam de modo estável.
\textbf{19.} Perto de $\pi$, com $\Theta = \pi + \epsilon$:
$\ddot\epsilon = [g/\ell - a^2\Omega^2/2\ell^2]\,\epsilon$; a
estabilidade exige o colchete negativo: $a^2\Omega^2 > 2g\ell$.
\textbf{20.} $\Omega_{\text{c}} = \sqrt{2g\ell}/a = \qty{140}{rad/s}$:
$f_{\text{c}} = \qty{22}{Hz}$; velocidade máxima $a\Omega_{\text{c}} =
\qty{2.8}{m/s}$, aceleração máxima $a\Omega_{\text{c}}^2 =
\qty{392}{m/s^2} \approx 40g$.
\textbf{21.} $\omega_{\text{lento}} = \sqrt{a^2\Omega^2/2\ell^2 -
g/\ell}$; em $\Omega = 2\Omega_{\text{c}}$, $a^2\Omega^2 = 8g\ell$, logo
$\omega_{\text{lento}} = \sqrt{3g/\ell} = \qty{8.6}{rad/s}$:
\qty{1.4}{Hz}, trinta vezes mais lento que a excitação de \qty{45}{Hz}.
\textbf{22.} $\xi_{\max} = (a/\ell)\sin\Theta$: inclinado $10^\circ$ em
relação à vertical, $\xi_{\max} = 0.05\sin170^\circ = \qty{8.7}{mrad}
\approx 0.5^\circ$ --- uma fotografia comum mostra uma vassoura parada
em pé.
\textbf{23.} O poço da posição em pé alcança os máximos que o ladeiam:
$|\Theta - \pi| < \arccos(2g\ell/a^2\Omega^2)$; em $\Omega =
2\Omega_{\text{c}}$, $\arccos\tfrac14 = 75^\circ$ --- um poço
notavelmente tolerante.
\textbf{24.} O malabarista usa realimentação: os olhos medem a
inclinação e a mão acelera lateralmente para anulá-la. A estabilização
de Kapitza é em malha aberta --- a excitação nunca sabe onde está o
pêndulo.
\textbf{25.} Sacudido a \qty{45}{Hz} com amplitude de \qty{2}{cm}, o
pêndulo de \qty{40}{cm} fica invertido num poço de $75^\circ$,
balançando a \qty{1.4}{Hz}. O mesmo potencial efetivo médio, obtido com
um campo elétrico quadrupolar oscilante em vez de um pivô sacudido,
confina íons isolados na armadilha de Paul --- o instrumento de
trabalho dos relógios atômicos e da computação quântica com íons
aprisionados.
\end{solution}
